\documentclass{ctexart}

% 中文支持由 ctexart 文档类自动提供

% 页面设置
\usepackage[top=2.54cm, bottom=2.54cm, left=3.18cm, right=3.18cm]{geometry}

% 数学公式
\usepackage{amsmath, amssymb, amsthm}

% 图形与表格
\usepackage{graphicx}
\usepackage{float}
\usepackage{booktabs}

% 超链接
\usepackage{hyperref}

% 量子态和算符
\usepackage{bm}
\newcommand{\ket}[1]{|#1\rangle}
\newcommand{\bra}[1]{\langle #1|}
\newcommand{\braket}[2]{\langle #1|#2\rangle}

% 标题设置
\title{Problem 1: Preliminaries and Warm-up}
\author{}
\date{\today}

\begin{document}

\maketitle

\section*{Problem 1: Preliminaries and Warm-up}

\subsection*{1.}
Show that $R_{\mathbf{n}}(\theta) = e^{-i\frac{\theta}{2}\mathbf{n}\cdot\boldsymbol{\sigma}}$ takes the form
\begin{equation}
    \cos(\theta/2)I - i\sin(\theta/2)(n_X X + n_Y Y + n_Z Z),
\end{equation}
where $\mathbf{n}$ is a unit 3D vector and $\boldsymbol{\sigma} = (X, Y, Z)$.

\medskip
\noindent\textbf{Solution:}
首先计算 $(\mathbf{n}\cdot\boldsymbol{\sigma})^2$。Pauli 矩阵满足对易关系
$\{X,Y\} = \{Y,Z\} = \{Z,X\} = 0$ 和 $X^2 = Y^2 = Z^2 = I$，因此
\begin{align*}
    (\mathbf{n}\cdot\boldsymbol{\sigma})^2
    &= (n_X X + n_Y Y + n_Z Z)^2 \\
    &= n_X^2 X^2 + n_Y^2 Y^2 + n_Z^2 Z^2
       + n_X n_Y (XY+YX) + n_Y n_Z (YZ+ZY) + n_Z n_X (ZX+XZ) \\
    &= (n_X^2 + n_Y^2 + n_Z^2) I \\
    &= I,
\end{align*}
最后一步利用了 $\mathbf{n}$ 是单位向量。

基于此，将指数函数作 Taylor 展开并按奇偶项分组：
\begin{align*}
    R_{\mathbf{n}}(\theta) = e^{-i\frac{\theta}{2}\mathbf{n}\cdot\boldsymbol{\sigma}}
    &= \sum_{k=0}^{\infty} \frac{1}{k!} \left(-i\frac{\theta}{2}\right)^k (\mathbf{n}\cdot\boldsymbol{\sigma})^k \\
    &= \sum_{m=0}^{\infty} \frac{(-i\frac{\theta}{2})^{2m}}{(2m)!} (\mathbf{n}\cdot\boldsymbol{\sigma})^{2m}
       + \sum_{m=0}^{\infty} \frac{(-i\frac{\theta}{2})^{2m+1}}{(2m+1)!} (\mathbf{n}\cdot\boldsymbol{\sigma})^{2m+1}.
\end{align*}
对于偶数项：$(\mathbf{n}\cdot\boldsymbol{\sigma})^{2m} = [(\mathbf{n}\cdot\boldsymbol{\sigma})^2]^m = I^m = I$；
对于奇数项：$(\mathbf{n}\cdot\boldsymbol{\sigma})^{2m+1} = (\mathbf{n}\cdot\boldsymbol{\sigma})^{2m}(\mathbf{n}\cdot\boldsymbol{\sigma}) = \mathbf{n}\cdot\boldsymbol{\sigma}$。

代入：
\begin{align*}
    R_{\mathbf{n}}(\theta)
    &= \sum_{m=0}^{\infty} \frac{(-1)^m (\theta/2)^{2m}}{(2m)!}\, I
       - i\,(\mathbf{n}\cdot\boldsymbol{\sigma}) \sum_{m=0}^{\infty} \frac{(-1)^m (\theta/2)^{2m+1}}{(2m+1)!} \\
    &= \cos(\theta/2)\,I \;-\; i\sin(\theta/2)\,(\mathbf{n}\cdot\boldsymbol{\sigma}) \\
    &= \cos(\theta/2)I - i\sin(\theta/2)(n_X X + n_Y Y + n_Z Z).
\end{align*}
这里我们利用了 $(-i)^{2m} = (-1)^m$ 和 $(-i)^{2m+1} = -i(-1)^m$。

\hfill $\square$

\subsection*{2.}
Show that any single-qubit unitary $U$ takes the form $e^{i\alpha}R_{\mathbf{n}}(\theta)$ for some $\alpha$, $\mathbf{n}$, and $\theta$.

\medskip
\noindent\textbf{Solution:}
任意 $2\times 2$ 酉矩阵 $U$ 满足 $U^\dagger U = I$ 且 $|\det U| = 1$。设 $\det U = e^{i\gamma}$，
则可提取全局相位：
\begin{equation*}
    U = e^{i\gamma/2} \, S, \qquad S \in \mathrm{SU}(2),\quad \det S = 1.
\end{equation*}
令 $\alpha = \gamma/2$。下面只需证明任意 $\mathrm{SU}(2)$ 矩阵 $S$ 均可写为 $R_{\mathbf{n}}(\theta)$ 的形式。

$\mathrm{SU}(2)$ 矩阵的最一般形式为
\begin{equation*}
    S = \begin{pmatrix}
        a & -b^* \\[2pt]
        b & a^*
    \end{pmatrix}, \qquad |a|^2 + |b|^2 = 1.
\end{equation*}
令 $a = a_r + i a_i$，$b = b_r + i b_i$，则约束为 $a_r^2 + a_i^2 + b_r^2 + b_i^2 = 1$。

另一方面，由第1问的结果：
\begin{equation*}
    R_{\mathbf{n}}(\theta)
    = \cos\frac{\theta}{2} I - i\sin\frac{\theta}{2}(n_X X + n_Y Y + n_Z Z)
    = \begin{pmatrix}
        \cos\frac{\theta}{2} - i n_Z \sin\frac{\theta}{2} & -i(n_X - i n_Y)\sin\frac{\theta}{2} \\[4pt]
        -i(n_X + i n_Y)\sin\frac{\theta}{2} & \cos\frac{\theta}{2} + i n_Z \sin\frac{\theta}{2}
    \end{pmatrix}.
\end{equation*}
对比 $S$ 的形式，取
\begin{equation*}
    a = \cos\frac{\theta}{2} - i n_Z \sin\frac{\theta}{2}, \qquad
    b = -i(n_X + i n_Y)\sin\frac{\theta}{2},
\end{equation*}
则有 $a_r = \cos\frac{\theta}{2}$, $a_i = -n_Z\sin\frac{\theta}{2}$, $b_r = n_Y\sin\frac{\theta}{2}$, $b_i = -n_X\sin\frac{\theta}{2}$。
由 $|\mathbf{n}|=1$ 可知 $|a|^2+|b|^2 = \cos^2\frac{\theta}{2} + \sin^2\frac{\theta}{2} = 1$，
即该映射覆盖了所有 $\mathrm{SU}(2)$ 矩阵。

具体地，给定额外的约束 $S^\dagger S = I$ 及 $\det S = 1$，可以直接解出参数：
\begin{align*}
    \theta &= 2\arccos(\mathrm{Re}(a)) = 2\arccos(a_r), \\
    n_X    &= -\frac{\mathrm{Im}(b)}{\sin(\theta/2)}, \quad
    n_Y    = \frac{\mathrm{Re}(b)}{\sin(\theta/2)}, \quad
    n_Z    = -\frac{\mathrm{Im}(a)}{\sin(\theta/2)}.
\end{align*}
由于 $a_r^2 + a_i^2 + b_r^2 + b_i^2 = 1$，$n_X^2+n_Y^2+n_Z^2 = 1$ 自然满足。

因此任意单量子比特酉变换 $U$ 都可写为 $U = e^{i\alpha}R_{\mathbf{n}}(\theta)$。

\hfill $\square$

\subsection*{3.}
Implement $R_{\mathbf{n}}(\theta)$ using $R_Y(\alpha)$ and $R_Z(\beta)$ for appropriate choices of $\alpha$ and $\beta$.

\medskip
\noindent\textbf{Solution:}
采用 Euler 角分解。将单位向量 $\mathbf{n}$ 写成球坐标形式：
\begin{equation*}
    \mathbf{n} = (\sin\beta\cos\alpha,\; \sin\beta\sin\alpha,\; \cos\beta),
\end{equation*}
其中 $\alpha \in [0, 2\pi)$，$\beta \in [0, \pi]$。

此时
\begin{equation*}
    \mathbf{n}\cdot\boldsymbol{\sigma}
    = \sin\beta\cos\alpha\, X + \sin\beta\sin\alpha\, Y + \cos\beta\, Z.
\end{equation*}

利用 Pauli 矩阵的相似变换：
\begin{align*}
    e^{-i\frac{\alpha}{2}Z} X e^{i\frac{\alpha}{2}Z} &= \cos\alpha\, X + \sin\alpha\, Y, \\
    e^{-i\frac{\beta}{2}Y} Z e^{i\frac{\beta}{2}Y}  &= \cos\beta\, Z + \sin\beta\, X,
\end{align*}
可以构造
\begin{align*}
    e^{-i\frac{\alpha}{2}Z} e^{-i\frac{\beta}{2}Y}\, Z\,
    e^{i\frac{\beta}{2}Y} e^{i\frac{\alpha}{2}Z}
    &= e^{-i\frac{\alpha}{2}Z} (\cos\beta\, Z + \sin\beta\, X) e^{i\frac{\alpha}{2}Z} \\
    &= \cos\beta\, Z + \sin\beta\cos\alpha\, X + \sin\beta\sin\alpha\, Y \\
    &= \mathbf{n}\cdot\boldsymbol{\sigma}.
\end{align*}

即 $\mathbf{n}\cdot\boldsymbol{\sigma} = R_Z(\alpha) R_Y(\beta)\, Z\, R_Y(-\beta) R_Z(-\alpha)$，其中
\begin{equation*}
    R_Y(\phi) = e^{-i\frac{\phi}{2}Y}, \qquad
    R_Z(\phi) = e^{-i\frac{\phi}{2}Z}.
\end{equation*}

对任意解析函数 $f$，有 $f(R\,A\,R^\dagger) = R\,f(A)\,R^\dagger$。令 $f(x) = e^{-i\frac{\theta}{2}x}$，则
\begin{align*}
    R_{\mathbf{n}}(\theta) = e^{-i\frac{\theta}{2}\mathbf{n}\cdot\boldsymbol{\sigma}}
    &= R_Z(\alpha) R_Y(\beta)\; e^{-i\frac{\theta}{2}Z}\; R_Y(-\beta) R_Z(-\alpha) \\
    &= R_Z(\alpha)\, R_Y(\beta)\, R_Z(\theta)\, R_Y(-\beta)\, R_Z(-\alpha).
\end{align*}

这就是最一般的分解式。特别地，当 $\mathbf{n}$ 位于 $XZ$ 平面时（即 $\alpha = 0$，$\mathbf{n} = (\sin\beta, 0, \cos\beta)$），
退化为更简洁的形式：
\begin{equation*}
    R_{\mathbf{n}}(\theta) = R_Y(\beta)\, R_Z(\theta)\, R_Y(-\beta).
\end{equation*}

因此 $R_{\mathbf{n}}(\theta)$ 总可以通过 $R_Y$ 和 $R_Z$ 门的组合来实现，具体参数由 $\mathbf{n}$ 的球坐标 $(\alpha, \beta)$ 和旋转角 $\theta$ 共同确定。

\hfill $\square$

\subsection*{4. State Preparation}
Starting with all qubits initialized in $\ket{0}$, prepare the following states using single-qubit gates,
two-qubit gates, ancillary qubits, classical feedback, and measurements. The resource cost
of the final circuit should be as small as possible. You may optimize for circuit size, circuit
depth, or both simultaneously.

\begin{enumerate}
    \item[(a)] Prepare an $n$-qubit quantum state $\frac{1}{\sqrt{2^n}} \sum_{x \in \{0,1\}^n} \ket{x}$;
    \item[(b)] Prepare an $n$-qubit state $\alpha\ket{000\cdots000} + \sqrt{1-\alpha^2}\ket{111\cdots111}$ for any $\alpha \in \mathbb{R}$;
    \item[(c)] Prepare the 4-qubit state $\frac{\sqrt{3}}{\sqrt{10}}\ket{0001} + \frac{2}{\sqrt{5}}\ket{0010} + \frac{\sqrt{6}}{4}\ket{0100} + \frac{\sqrt{6}}{4}\ket{1000}$;
    \item[(d)] (Optional) Prepare the following state:
    \begin{align*}
        &\alpha(\ket{1000000}+\ket{0010101}+\ket{1110011}+\ket{0100110} \\
        &\quad +\ket{1001111}+\ket{0011010}+\ket{1111100}+\ket{0101001}) \\
        &+\beta(\ket{0111111}+\ket{1101010}+\ket{0001100}+\ket{1011001} \\
        &\quad +\ket{0110000}+\ket{1100101}+\ket{0000011}+\ket{1010110})
    \end{align*}
    \item[(e)] (Optional) Use TensorCircuit to verify your answers.
\end{enumerate}

\subsection*{5. Gate Construction}

\begin{enumerate}
    \item[(a)] Use CNOT gates to implement the SWAP operation: $\text{SWAP}\ket{\psi}\ket{\phi} = \ket{\phi}\ket{\psi}$.
    \item[(b)] Consider four qubits $q_0, q_1, q_2$ and $q_3$. Implement two CNOT gates between $q_0$ and $q_2$, and $q_1$ and $q_3$ respectively, under the constraint that only CNOT gates between $q_i$ and $q_{i+1}$ are available. Try to use as few gates as possible.
    \item[(c)] Use CNOT and $R_Z$ gates to implement the two-qubit gate $R_{ZZ}$, where
    \begin{equation*}
        R_Z(\theta) = e^{-i\theta/2 Z}, \qquad R_{ZZ}(\theta) = e^{-i\theta/2 Z\otimes Z}.
    \end{equation*}
    \item[(d)] Construct a Controlled-Phase gate
    \begin{equation*}
        \begin{pmatrix}
            1 & 0 & 0 & 0 \\
            0 & 1 & 0 & 0 \\
            0 & 0 & 1 & 0 \\
            0 & 0 & 0 & e^{i\theta}
        \end{pmatrix}
    \end{equation*}
    using $R_{ZZ}$ and single-qubit gates.
    \item[(e)] For an arbitrary $n$-qubit Pauli operator $P$, construct a circuit to implement $e^{-i\theta P}$.
    \item[(f)] Implement the following specific $\mathrm{SU}(4)$ gate:
    \begin{equation*}
        \begin{bmatrix}
            \frac{1}{2} & \frac{1}{\sqrt{2}} & \frac{1}{2} & 0 \\[6pt]
            \frac{1}{\sqrt{2}} & -\frac{1}{2} & \frac{1}{\sqrt{2}} & 0 \\[6pt]
            \frac{\sqrt{3}}{2\sqrt{2}} & 0 & -\frac{\sqrt{3}}{2\sqrt{2}} & \frac{\sqrt{3}}{2} \\[6pt]
            \frac{\sqrt{3}}{2\sqrt{2}} & 0 & -\frac{\sqrt{3}}{2\sqrt{2}} & -\frac{1}{2}
        \end{bmatrix}
    \end{equation*}
    using single-qubit pulses and CZ gates.
    \item[(g)] (Optional) Use TensorCircuit to verify your answer above.
\end{enumerate}

\subsection*{6. Measurement}

\begin{enumerate}
    \item[(a)] Construct a quantum circuit to measure the two-qubit operators $Z_1 Z_2$ and $X_1 X_2$ on $\ket{\psi}$ and project $\ket{\psi}$ onto the joint $+1$ eigenstate of $Z_1 Z_2$ and $X_1 X_2$.
    \item[(b)] Give a general method to construct a quantum circuit to perform projective measurement of any Pauli operator $P$.
    \item[(c)] Prepare the single-qubit quantum state $\alpha\ket{0} + \sqrt{1-\alpha^2}\ket{1}$ from an $n$-qubit quantum state $\alpha\ket{000\cdots000} + \sqrt{1-\alpha^2}\ket{111\cdots111}$ by measurement.
    \item[(d)] (Optional) A 5-qubit fanout gate $F_4$ acting on qubits $q_0, q_1, q_2, q_3, q_4$ is defined as follows
    \begin{equation*}
        F_4 \ket{x_0}\ket{x_1, x_2, x_3, x_4} = \ket{x_0}\ket{x_1 \oplus x_0, x_2 \oplus x_0, x_3 \oplus x_0, x_4 \oplus x_0},
    \end{equation*}
    where $x_0, x_1, \ldots, x_4 \in \{0,1\}$. Two-qubit gates may act only on qubits $(q_j, q_{j+1})$ for $0 \leq j \leq 3$. Construct the quantum circuit of $F_4$ to minimize the circuit depth and size. (Hint: use ancillary qubits and measurements.)
    \item[(e)] (Optional) Use CNOTs and Hadamards to compute the inner product $|\braket{\psi|\phi}|$ for arbitrary states $\ket{\psi}$ and $\ket{\phi}$.
\end{enumerate}

\subsection*{7. Reversible Circuit}
(Optional) Construct a reversible classical circuit using Toffoli and NOT gates to evaluate the formula:
\begin{align*}
    S_3(x) = &(\neg x_1 \lor \neg x_2) \land (x_1 \lor \neg x_2) \land (\neg x_1 \lor x_2 \lor \neg x_3) \\
             &\land (\neg x_1 \lor x_2 \lor x_3) \land (x_1 \lor x_2 \lor x_3),
\end{align*}
where each $x_i$ represents a Boolean value, and $\neg$, $\lor$, and $\land$ denote logical NOT, OR, and AND, respectively. Try to: (1) use as few ancillary qubits as possible, and (2) minimize the product of the number of gates and the total number of qubits.

\subsection*{8. Grover Search}
(Optional) Using the results from the previous problem, find a configuration of $x$ satisfying $S_3(x) = \text{True}$ with Grover's search algorithm. Implement your algorithm using TensorCircuit and verify the solution.

\end{document}
